\documentclass[11pt]{article}
\usepackage{../../styles/activity}

\usepackage{xr}
\externaldocument{0-MR}

\lhead{}
%\chead{\textbf{\Large{\hspace{0pt}Beginning Activities for Section~5.5}}\\\hspace{0pt}\emph{Mathematical Reasoning: Writing and Proof}}
\bahead{5.5}
\rhead{}
\lfoot{}
\rfoot{}
\cfoot{\hspace{0pt}\scalebox{0.4}{\includegraphics{cc-by-nc-sa.eps}}}

\begin{document}
\input table

\subsection*{Beginning Activity 1 (The Union and Intersection of a Family of Sets)}

\begin{enumerate}
\item \begin{multicols}{2}
$A \cup B \cup C = \left\{1, 2, 3, 4, 5, 6, 7 \right\}$ \\
$B \cup C \cup D = \left\{2, 3, 4, 5, 6, 7, 8 \right\}$ \\
$A \cap B \cap C = \left\{3, 4, 5 \right\}$ \\
$B \cap C \cap D = \left\{4, 5, 6 \right\}$ 
\end{multicols}

\item 
\begin{enumerate}
\item $( A \cup B \cup C  ) \cup D = \left\{1, 2, 3, 4, 5, 6, 7, 8 \right\}$
\item $A \cup ( B \cup C  \cup D ) = \left\{1, 2, 3, 4, 5, 6, 7, 8 \right\}$
\item $A \cup ( B \cup C )   \cup D = \left\{1, 2, 3, 4, 5, 6, 7, 8 \right\}$
\item $( A \cup B ) \cup ( C \cup D ) = \left\{1, 2, 3, 4, 5, 6, 7, 8 \right\}$
\item $( A \cap B \cap C  ) \cap D = \left\{4, 5 \right\}$
\item $A \cap ( B \cap C  \cap D ) = \left\{4, 5 \right\}$
\item $A \cap ( B \cap C )   \cap D = \left\{4, 5 \right\}$
\item $( A \cap B ) \cap ( C \cap D ) = \left\{4, 5 \right\}$
\end{enumerate}

\item It appears that the placement of the parentheses in the union or intersection of these four sets does not make a difference.  This means that we can define $A \cup B \cup C \cup D$ and $A \cap B \cap C \cap D$ as $( A \cup B \cup C  ) \cup D$ and 
$( A \cap B \cap C  ) \cap D$, respectively.

\item Since $\mathscr{A} = \left\{A, B, C , D \right\}$,
\begin{itemize}
\item $\bigcup\limits_{X \in \mathscr{A}}^{}X$ consists of all the elements that are in one of the four sets $A$, $B$, $C$, and $D$.  This is the same set as  $A \cup B \cup C \cup D$.

\item $\bigcap\limits_{X \in \mathscr{A}}^{}X$ consists of all the elements that are in all of the four sets $A$, $B$, $C$, and $D$.  This is the same set as  $A \cap B \cap C \cap D$.
\end{itemize}  
So,
\[
\bigcup_{X \in \mathscr{A}}^{}X = \left\{1, 2, 3, 4, 5, 6, 7, 8 \right\} \qquad \text{and} \qquad
\bigcap_{X \in \mathscr{A}}^{}X = \left\{4, 5 \right\}.
\]

\item $\bigcup\limits_{X \in \mathscr{B}}^{}X = \left\{1, 2, 3, 4, 5, 6, 7, 8, 9, 10 \right\}$ \qquad and \qquad $\bigcap_{X \in \mathscr{B}}^{}X = \emptyset$

\item Recall the the universal set is $\N$.  Therefore,
\[
\left( \bigcup_{X \in \mathscr{A}}^{}X \right)^c = \left\{9, 10, 11, 12, 13, \ldots \right\}.
\]
In addition,
\begin{multicols}{2}
$A^c = \left\{6, 7, 8, 9, \ldots \right\}$ \\
$B^c = \left\{7, 8, 9, 10, \ldots \right\}$ \\
$C^c = \left\{8, 9, 10, 11, \ldots \right\}$ \\
$D^c = \left\{9, 10, 11, 12, \ldots \right\}$. 
\end{multicols}
Therefore,
\[
\bigcap_{X \in \mathscr{A}}^{}X^c = \left\{9, 10, 11, 12, \ldots \right\}.
\]
\end{enumerate}
\hbreak



\subsection*{Beginning Activity 2 (An Indexed Family of Sets)}
\begin{enumerate}
\item $\bigcup\limits_{j=1}^{4}C_j = \left\{ 1, 2, 3, 4, 5, 6, 7, 8 \right\}$ \qquad and  \qquad
$\bigcap\limits_{j=1}^{4}C_j = \left\{ 4, 5 \right\}$

\item \begin{enumerate}
\item $\bigcup\limits_{j=1}^{6}C_j = \left\{ 1, 2, 3, 4, 5, 6, 7, 8, 9, 10 \right\}$ and  
$\bigcap\limits_{j=1}^{6}C_j = \emptyset$

\item $\bigcup\limits_{j=1}^{8}C_j = \left\{ 1, 2, 3, 4, 5, 6, 7, 8, 9, 10, 11, 12 \right\}$ and  $\bigcap\limits_{j=1}^{8}C_j = \emptyset$

\item $\bigcup\limits_{j=4}^{8}C_j = \left\{ 4, 5, 6, 7, 8, 9, 10, 11, 12 \right\}$ and  
$\bigcap\limits_{j=1}^{8}C_j = \left\{ 8 \right\}$


\item $( \bigcap\limits_{j=1}^{4}C_j )^c = \left\{1, 2, 3, 6, 7, 8, 9, 10, \ldots \right\}$ and \\
$\bigcup\limits_{j=1}^{4}C_j^c = \left\{1, 2, 3, 6, 7, 8, 9, 10, \ldots \right\}$

\end{enumerate}
\end{enumerate}
\hbreak

\end{document}
